\documentclass[11pt]{article}
\usepackage[utf8]{inputenc}
\usepackage[T1]{fontenc}
\usepackage[spanish,es-nodecimaldot]{babel}
\usepackage[a4paper,margin=2.5cm]{geometry}
\usepackage{amsmath,amssymb,booktabs,graphicx,float,hyperref}
\newcommand{\figroot}{../../outputs/toi_ati_case_anatomy/figures_pdf}
\newcommand{\safeincludegraphics}[2][]{%
  \IfFileExists{#2}{\includegraphics[#1]{#2}}{\fbox{\begin{minipage}{0.86\textwidth}\centering Figura no encontrada: \texttt{#2}\end{minipage}}}%
}
\title{Anatomia de casos TOI/ATI: regiones Mapper y planetas ancla}
\author{Proyecto Mapper/TDA de exoplanetas}
\date{}
\begin{document}
\maketitle

\section{Introduccion}
Este reporte abre el ranking global TOI/ATI para explicar por que ciertas regiones Mapper y planetas ancla fueron priorizados. El objetivo no es afirmar descubrimientos de planetas, sino auditar la anatomia de los indices.

\section{Que problema resuelve la anatomia TOI/ATI}
La salida principal es una anatomia de los casos top: que factor empuja cada indice, que factor lo limita y que tan estable parece el deficit relativo al cambiar de radio.

\section{Datos usados}
El modulo consume tablas generadas por \texttt{topological\_incompleteness\_index} y, cuando existen, tablas contextuales de sombra observacional y estudios locales.

\section{Descomposicion de TOI}
\[
\mathrm{TOI}(v)=\mathrm{Shadow}(v)(1-I_{R^3}(v))C_{\mathrm{phys}}(v)S_{\mathrm{net}}(v).
\]
TOI debe leerse como una prioridad regional que combina frontera observacional, baja imputacion, continuidad fisica con vecinos y soporte de red.

\section{Descomposicion de ATI}
\[
\mathrm{ATI}(p^*)=\mathrm{TOI}(v)\Delta_{\mathrm{rel,best}}(p^*)(1-I_{R^3}(p^*))A(p^*).
\]
ATI debe leerse como una prioridad local de inspeccion. El termino $\Delta_{\mathrm{rel,best}}$ debe revisarse junto con los radios individuales.

\section{Top regiones}
\begin{figure}[H]\centering
\safeincludegraphics[width=.9\textwidth]{\figroot/top_regions_toi_case_anatomy.pdf}
\caption{Top regiones por TOI.}
\end{figure}

\begin{figure}[H]\centering
\safeincludegraphics[width=.9\textwidth]{\figroot/toi_factor_decomposition.pdf}
\caption{Descomposicion visual de factores TOI.}
\end{figure}

\section{Top planetas ancla}
\begin{figure}[H]\centering
\safeincludegraphics[width=.9\textwidth]{\figroot/top_anchors_ati_case_anatomy.pdf}
\caption{Top planetas ancla por ATI.}
\end{figure}

\begin{figure}[H]\centering
\safeincludegraphics[width=.9\textwidth]{\figroot/ati_factor_decomposition.pdf}
\caption{Descomposicion visual de factores ATI.}
\end{figure}

\section{Tablas de deficit por radio para las anclas principales}
No basta con usar $\Delta_{\mathrm{rel,best}}$ porque ese maximo puede inflar la lectura de un caso local. Por eso se reportan los tres radios: $r_{\mathrm{kNN}}$, $r_{\mathrm{node\_median}}$ y $r_{\mathrm{node\_q75}}$. Cada tabla muestra $N_{\mathrm{obs}}$, la referencia local $N_{\mathrm{exp}}$, el deficit absoluto $\Delta N$ y el deficit relativo $\Delta_{\mathrm{rel}}$.

La diferencia clave es que $\Delta N$ depende de la escala de conteo, mientras que $\Delta_{\mathrm{rel}}$ normaliza por la referencia esperada. Si el deficit permanece positivo en los tres radios, el caso es mas estable. Si solo aparece en una escala, debe presentarse como exploratorio.

\subsection{Deficit por radio para \texttt{HIP 97166 c} en \texttt{cube12\_cluster0}}
\begin{table}[H]
\centering
\small
\caption{Conteos observados, esperados y deficit local para el planeta ancla \texttt{HIP 97166 c}.}
\begin{tabular}{lrrrr}
\toprule
Radio & $N_{\mathrm{obs}}$ & $N_{\mathrm{exp}}$ & $\Delta N$ & $\Delta_{\mathrm{rel}}$ \\
\midrule
$r_{\mathrm{kNN}}$ & 10.00 & 11.00 & 1.00 & 0.091 \\
$r_{\mathrm{node\_median}}$ & 1617.00 & 1144.00 & -473.00 & -0.413 \\
$r_{\mathrm{node\_q75}}$ & 2550.00 & 1914.00 & -636.00 & -0.332 \\
\bottomrule
\end{tabular}
\end{table}
El deficit depende de la escala local elegida. Esto sugiere que el ancla es util para priorizacion, pero la evidencia es sensible al radio y debe considerarse exploratoria.

\subsection{Deficit por radio para \texttt{HIP 90988 b} en \texttt{cube17\_cluster2}}
\begin{table}[H]
\centering
\small
\caption{Conteos observados, esperados y deficit local para el planeta ancla \texttt{HIP 90988 b}.}
\begin{tabular}{lrrrr}
\toprule
Radio & $N_{\mathrm{obs}}$ & $N_{\mathrm{exp}}$ & $\Delta N$ & $\Delta_{\mathrm{rel}}$ \\
\midrule
$r_{\mathrm{kNN}}$ & 10.00 & 11.00 & 1.00 & 0.091 \\
$r_{\mathrm{node\_median}}$ & 8.00 & 9.00 & 1.00 & 0.111 \\
$r_{\mathrm{node\_q75}}$ & 62.00 & 63.00 & 1.00 & 0.016 \\
\bottomrule
\end{tabular}
\end{table}
El deficit aparece en las tres escalas locales, por lo que la evidencia topologica de submuestreo es mas estable. Aun asi, no debe interpretarse como conteo absoluto de objetos reales ausentes.

\subsection{Deficit por radio para \texttt{HD 42012 b} en \texttt{cube19\_cluster5}}
\begin{table}[H]
\centering
\small
\caption{Conteos observados, esperados y deficit local para el planeta ancla \texttt{HD 42012 b}.}
\begin{tabular}{lrrrr}
\toprule
Radio & $N_{\mathrm{obs}}$ & $N_{\mathrm{exp}}$ & $\Delta N$ & $\Delta_{\mathrm{rel}}$ \\
\midrule
$r_{\mathrm{kNN}}$ & 10.00 & 11.00 & 1.00 & 0.091 \\
$r_{\mathrm{node\_median}}$ & 120.00 & 121.00 & 1.00 & 0.008 \\
$r_{\mathrm{node\_q75}}$ & 293.00 & 294.00 & 1.00 & 0.003 \\
\bottomrule
\end{tabular}
\end{table}
El deficit aparece en las tres escalas locales, por lo que la evidencia topologica de submuestreo es mas estable. Aun asi, no debe interpretarse como conteo absoluto de objetos reales ausentes.

\subsection{Deficit por radio para \texttt{HD 42012 b} en \texttt{cube26\_cluster6}}
\begin{table}[H]
\centering
\small
\caption{Conteos observados, esperados y deficit local para el planeta ancla \texttt{HD 42012 b}.}
\begin{tabular}{lrrrr}
\toprule
Radio & $N_{\mathrm{obs}}$ & $N_{\mathrm{exp}}$ & $\Delta N$ & $\Delta_{\mathrm{rel}}$ \\
\midrule
$r_{\mathrm{kNN}}$ & 10.00 & 11.00 & 1.00 & 0.091 \\
$r_{\mathrm{node\_median}}$ & 120.00 & 121.00 & 1.00 & 0.008 \\
$r_{\mathrm{node\_q75}}$ & 293.00 & 294.00 & 1.00 & 0.003 \\
\bottomrule
\end{tabular}
\end{table}
El deficit aparece en las tres escalas locales, por lo que la evidencia topologica de submuestreo es mas estable. Aun asi, no debe interpretarse como conteo absoluto de objetos reales ausentes.

\subsection{Deficit por radio para \texttt{HD 4313 b} en \texttt{cube17\_cluster10}}
\begin{table}[H]
\centering
\small
\caption{Conteos observados, esperados y deficit local para el planeta ancla \texttt{HD 4313 b}.}
\begin{tabular}{lrrrr}
\toprule
Radio & $N_{\mathrm{obs}}$ & $N_{\mathrm{exp}}$ & $\Delta N$ & $\Delta_{\mathrm{rel}}$ \\
\midrule
$r_{\mathrm{kNN}}$ & 10.00 & 11.00 & 1.00 & 0.091 \\
$r_{\mathrm{node\_median}}$ & 5.00 & 6.00 & 1.00 & 0.167 \\
$r_{\mathrm{node\_q75}}$ & 11.00 & 12.00 & 1.00 & 0.083 \\
\bottomrule
\end{tabular}
\end{table}
El deficit aparece en las tres escalas locales, por lo que la evidencia topologica de submuestreo es mas estable. Aun asi, no debe interpretarse como conteo absoluto de objetos reales ausentes.



\section{Deficit por radio}
\begin{figure}[H]\centering
\safeincludegraphics[width=.8\textwidth]{\figroot/deficit_by_radius_summary.pdf}
\caption{Resumen del deficit relativo por radio.}
\end{figure}

\section{Tres casos finales para exposicion}
El caso regional no afirma planetas faltantes; prioriza una zona Mapper. El caso ancla no afirma un objeto ausente; prioriza una vecindad local. El caso repetido no es duplicacion erronea; puede ser una senal de transicion topologica por solapamiento de cubiertas.

Si la tabla por radio muestra que $\Delta_{\mathrm{rel}}$ es alto solo en una escala, el caso debe presentarse como exploratorio. Si el deficit se mantiene positivo en los tres radios, el caso es mas estable.

\begin{table}[H]
\centering
\small
\caption{Tres casos finales para exposicion.}
\begin{tabular}{p{2.5cm}p{2.8cm}p{2.4cm}p{3.8cm}p{3.8cm}}
\toprule
Caso & Objeto & Nodo & Motivo de seleccion & Cautela principal \\
\midrule
top\_toi\_region & cube12\_cluster0 & cube12\_cluster0 & Presentar como la region Mapper con mayor prioridad regional. Explicar que gana por la combinacion de sombra observacional, baja imputacion, continuidad fisica y soporte de red. & El caso regional no afirma objetos ausentes; prioriza una zona Mapper para inspeccion observacional. \\
top\_ati\_anchor & HIP 97166 c & cube12\_cluster0 & Presentar como el planeta ancla mas prioritario para inspeccion local. Explicar que combina una region TOI alta con deficit local y baja imputacion. & El caso ancla no afirma un objeto ausente; prioriza una vecindad local en R^3. \\
repeated\_anchor\_multi\_node & HD 42012 b & cube26\_cluster6 & Presentar como un caso de transicion Mapper. El planeta aparece en varios nodos por el solapamiento de la cubierta, lo que puede indicar que vive en una frontera topologica entre vecindarios. & El caso repetido no es duplicacion erronea; puede reflejar solapamiento de cubiertas y transicion topologica. \\
\bottomrule
\end{tabular}
\end{table}

\begin{figure}[H]\centering
\safeincludegraphics[width=.95\textwidth]{\figroot/final_presentation_cases_summary.pdf}
\caption{Resumen de los tres casos finales para exposicion.}
\end{figure}

\section{Discusion}
Una region con TOI alto puede ganar por sombra fuerte, baja imputacion, continuidad fisica o soporte de red. Un ancla con ATI alto puede ganar por pertenecer a una region TOI alta, por mostrar deficit local, por baja imputacion o por ser representativa del nodo.

\section{Limitaciones}
Este modulo no introduce una funcion de completitud instrumental, no estima cantidades absolutas de objetos reales ausentes y no convierte por si solo un caso priorizado en una conclusion observacional cerrada.

\section{Conclusion}
Este modulo no calcula otro indice nuevo; abre los rankings TOI/ATI para explicar por que ciertas regiones y planetas ancla fueron priorizados. La salida principal es una anatomia de los casos top, no una afirmacion de objetos ausentes confirmados.
\end{document}
